\section{Sharp zero multiplicities for symmetric Galois groups}
\label{sec:reflection}

An incoming report considers the splitting field of a polynomial with
symmetric Galois group. It asks whether the order of an off-line zero
can be recovered with the optimal loss from one much smaller set of
Frobenius classes. Its symmetric-group estimate is
$2(n-H_n)$, where $H_n=\sum_{j=1}^n1/j$.
It gives $n-1$ as the best possible lower bound for a proof from local
character positivity. The following result attains that bound for
every $3\le n\le 12$ and also controls zeros on the critical line. The precise
antecedent is \emph{Reflection detectors for nonabelian Riemann
hypotheses}, the symmetric-group corollary and the question in its final section,
in the \href{https://github.com/VladimirReshetnikov/ProveIt/blob/5cc267af75a635199ac5d93a7d49cbce521daefc/docs/incoming/grh_reflection_transfer_2026-11-25.zip}{pinned reflection report}.

\paragraph{Theorem.}
Let $L/K$ be a finite Galois extension of number fields with group
$G=S_n$, $3\le n\le 12$. Fix a transposition $t$, set
$E=L^{\langle t\rangle}$ and $F=L^{A_n}$, and let $\varepsilon$
be the sign character. Then the meromorphic function
\[
 \mathcal Q(s)=\frac{\zeta_E(s)^{\,n-1}L(s,\varepsilon,L/K)}
                         {\zeta_L(s)}
 =\frac{\zeta_E(s)^{\,n-1}\zeta_F(s)}{\zeta_L(s)\zeta_K(s)}
\]
is holomorphic on $\mathbb C\setminus\{1\}$, with a pole of exact order
$n-2$ at~1. In particular, for $K=\mathbb Q$ and $1/2<\Re\rho<1$,
Dirichlet GRH gives
\[
 \operatorname{ord}_\rho\zeta_L\le (n-1)\operatorname{ord}_\rho\zeta_E .
\]
For example, the constants are $4$ and $5$ for $S_5$ and $S_6$. Each constant $n-1$ is optimal among deductions
using positivity against all characters induced from one-dimensional
subgroup characters and vanishing of the two global linear coefficients.

\paragraph{Arithmetic input.}
For $s\ne1$ let $\Theta=\sum_\chi (\operatorname{ord}_{s}L(\cdot,\chi,L/K))\chi$,
the sum over irreducible characters of $G$.
Brauer induction gives meromorphy, so this virtual character exists.
For every $H\le G$ and every linear character $\eta$ of $H$, Artin
induction formalism and Hecke holomorphy away from 1 give
$\langle\Theta|_H,\eta\rangle\ge0$. The same holds for
$\Psi=\varepsilon\Theta$ since $\varepsilon|_H\eta$ is linear.
The regular and permutation characters give
\[
 m:=\Psi(1)=\operatorname{ord}_s\zeta_L,\qquad
 a:=\tfrac12(m-\Psi(t))=\operatorname{ord}_s\zeta_E,\qquad
 b:=\langle\Psi,1\rangle=\operatorname{ord}_s L(\cdot,\varepsilon,L/K).
\]
These are the classical Heilbronn-character identities.
See H.~Heilbronn, \href{https://doi.org/10.4153/CJM-1973-090-3}{\emph{On real
zeros of Dedekind $\zeta$-functions}}, Canad. J. Math. 25 (1973), 870--873,
and R.~Foote, H.~Ginsberg and V.~K.~Murty,
\href{https://doi.org/10.1090/BULL/1492}{\emph{Heilbronn characters}}.

\paragraph{Finite proof.}
For a conjugacy type $T$, let $V_T$ be the value of $\Psi$ on $T$
and $c_T$ the size of the class. Write $\delta_T$ for its coordinate
unit vector. Given a subgroup $H$ and a surjection $e:H\to\mathbb Z/d$,
put
\[
 B_{H,e}(T)=\sum_{h\in H\cap T}\frac1{\varphi(d)}
       \sum_{u\in(\mathbb Z/d)^\times}
                   \exp(2\pi i u e(h)/d).
\]
The row $B_{H,e}$ is real and $B_{H,e}\cdot V\ge0$ by the preceding
local positivity. The sum over $u$ includes the complex conjugate of
each character. We claim the following exact vector identity in each certified degree. The table gives the positive weights for the two smallest nonsolvable groups.
\[
 \frac{n-3}{n-1}\delta_{1}-\delta_t+
             \frac{2}{(n-1)n!}(c_T)_T
       =\sum_{(H,e)} w_{H,e} B_{H,e}.                         \tag{R}
\]
Pairing with $V$ gives $2(a+(b-m)/(n-1))\ge0$, i.e.
$m\le(n-1)a+b$ at every $s\ne1$. This proves the holomorphy.
At $s=1$, each Dedekind zeta has a simple pole with nonzero residue,
which gives the stated exact order.

The notation $g_j$ in the following table means that the displayed
permutation maps to $j\bmod d$. Each row generates its indicated
subgroup and character. For example, the first row contains the identity,
two transpositions with character value $-1$, and one double
transposition with value $+1$.

\begin{center}\small
\begin{tabular}{r r r l}
\hline
$n$ & $w_{H,e}$ & $d$ & Generators with their images\\ \hline
5 & $3/8$&2 & $(45)_1,\ (23)_1$\\
5 & $1/16$&2 & $(45)_1,\ ((24)(35))_1$\\
5 & $1/24$&2 & $(45)_1,\ ((12)(34))_1$\\
5 & $1/40$&2 & $((23)(45))_1,\ ((12)(34))_1$\\ \hline
6 & $7/30$&2 & $(56)_1,\ ((23)(45))_1$\\
6 & $1/10$&3 & $(456)_2,\ ((12)(34)(56))_0$\\
6 & $1/45$&6 & $(456)_4,\ ((14)(25)(36))_3$\\
6 & $1/50$&2 & $((34)(56))_1,\ ((23)(45))_1$\\
6 & $1/5$&2 & $(56)_1,\ (34)_1,\ (12)_1$\\
6 & $1/40$&2 & $(56)_0,\ (34)_1,\ ((13)(24))_1$\\ \hline
\end{tabular}
\end{center}

For $S_5$ one can check (R) directly. In the type order
$1,2,3,22,4,32,5$ (fixed points suppressed), the four rows are
\[
\begin{pmatrix}
1&-2&0&1&0&0&0\\
1&-2&0&-1&2&0&0\\
1&-2&2&-3&0&2&0\\
1&0&0&-5&0&0&4
\end{pmatrix}.
\]
The class sizes are $(1,10,20,15,30,20,24)$.
To reproduce the checks for every degree, start with the identity
permutation at image zero and multiply by each generator, adding its
image modulo $d$, until the list stops growing. Each occurrence of an
already listed permutation has the same image. Sum the real cyclotomic
traces in each conjugacy class and check (R) entry by entry. The
accompanying \texttt{check\_reflection.py} performs these operations in
exact rational arithmetic, using the completely specified inputs in
\texttt{certificates.json}. The optimality witness is the usual standard
character $\Psi(g)=\#\operatorname{Fix}(g)-1$. For every $n\ge3$ it is
orthogonal to both linear characters, it satisfies every subgroup
positivity inequality, and it has $(m,a,b)=(n-1,1,0)$. This witnesses
sharpness of the group-theoretic certificate. Existence of an actual
off-line zero realizing equality is a separate question.

For $n=3$ the positive row is the sign character itself, weighted
by $1/6$ and aggregated by type. The accompanying data for each
remaining degree contain a list of generators and a cyclic character
on each subgroup in (R), as in the table. The number of columns in
(R) is the number of partitions of $n$, and its class sizes follow
from $c_T=n!/\prod_{j} j^{m_j}m_j!$ for a type with $m_j$ cycles
of length $j$. The checker enumerates the generated subgroup, with two class-count
shortcuts for large rows. A full $S_k$ is recognized by a $k$-cycle
and an adjacent transposition; its row uses the class-size formula and
the indicated sign power. An additional generator exchanging two
disjoint equal blocks gives $(S_k\times S_k)\rtimes C_2$. In the
base subgroup the cycle types concatenate. In the exchange coset the
cycles have twice the lengths of the cycles of the product of the
two component permutations, and there are $k!$ pairs with each
given product. The linear character is the product of its values on
the two components, times its value on the exchange in that coset.
The checker matches the supplied generators and their character values
before applying either identity. No numerical optimization enters its
acceptance test.

\begin{center}\small
\begin{tabular}{c|rrrrrrrrrr}
$n$&3&4&5&6&7&8&9&10&11&12\\
conjugacy classes&3&5&7&11&15&22&30&42&56&77\\
positive subgroup rows&1&3&4&6&10&18&26&38&48&71
\end{tabular}
\end{center}

\paragraph{The character behind the calculation.}
Write $\sigma$ for the nontrivial character of $\langle t\rangle$.
The character whose pairing with $\Psi$ is
$(n-1)a-m+b$ equals
\[
 U_n=(n-1)\operatorname{Ind}_{\langle t\rangle}^{S_n}\sigma
                   -\operatorname{Reg}_{S_n}+1.
\]
It is a genuine character for every $n$. Indeed, take the Cayley graph
of $S_n$ for the $n-1$ adjacent transpositions, with left multiplication
on vertices. Its vertex module is regular and its oriented edge module
has one summand $\operatorname{Ind}\sigma$ for each generator, because
the edge stabilizer reverses the edge. The graph is connected, so the
cellular chain sequence gives $U_n=H_1$ at the level of characters.
A proof that $U_n$ belongs to the positive rational cone of characters
induced from linear subgroup characters for all $n$ would give the
multiplicity theorem in all degrees. Genuine character status alone
gives no such decomposition. Formula (R) supplies it in the stated
degrees by Frobenius reciprocity.

There is an obstruction to deducing the required decomposition from graph
homology in general. Let $A_5$ fix two vertices and permute five parallel
edges in its natural action. Its first homology is the irreducible
four-dimensional augmentation representation $U$ (the action is doubly
transitive). If $U$ were a positive rational sum of monomial characters,
each summand would be a multiple $kU$. Frobenius reciprocity gives
$1\le k\le4$, so the inducing subgroup would have index $4k$.
Indices $8$ and $16$ do not divide $60$, and an index-four subgroup would
give a nontrivial map of the simple group $A_5$ into $S_4$. In the
remaining case the subgroup has order five and $k=3$. The restriction of
$U$ to this subgroup is its regular representation minus its trivial
character, so every linear character occurs at most once. This rules
out $k=3$. Hence the general-degree problem needs more about the
adjacent-transposition graph than connectedness and a group action.

\paragraph{What the improvement detects.}
For the $S_5$ splitting field of $X^5-X-1$ in the prior report, an
off-line zero of multiplicity $m$ would have multiplicity at least
$\lceil m/4\rceil$ in the degree-60 transposition fixed field.
The former constant was $163/30$. For $S_6$ the corresponding degree
is $360$ and the former constant was $71/10$.
Artin--Hecke formalism gives the new meromorphic quotient over
every base field, independently of GRH. Over $\mathbb Q$, the Lean
Dirichlet result removes its only possible linear-character correction
to the right of the critical line. In the previous report's proposed
one-sided prime estimate, this identifies exactly how much multiplicity
the selected Frobenius classes must retain. An estimate for their prime
sum is still required to deduce RH for the full splitting field.
